\documentclass[10pt,a4paper]{amsart}
\usepackage[margin=19mm]{geometry}
\usepackage{amsmath,amssymb,mathtools}
\usepackage[scr=rsfs,cal=euler]{mathalfa}
\usepackage{xfrac}
\usepackage[unicode,hidelinks,bookmarks=false]{hyperref}
\newcommand{\ie}{i.e.\ }
\newcommand{\cf}{cf.\ }
\newcommand{\resp}{respectively\ }
\newcommand{\Subsection}{\subsection}
\providecommand{\nobreakdash}{\nobreak\hbox{-}}
\setlength{\emergencystretch}{2em}
\multlinegap0pt
\usepackage{mathtools}
\usepackage{amssymb}
\usepackage{xcolor}
\usepackage{tikz}
\newtheorem{proposition}{Proposition}
\newcommand{\Li}{\mathrm{Li}}
\newcommand{\Z}{\mathbb{Z}}
\newcommand{\Zp}{\mathbb{Z}_p}
\title{Multifractal analysis of power means for the Schneider map on $p\mathbb{Z}_{p}$}
\author{Matias Alvarado, Nicol\'as Ar\'evalo-Hurtado}
\date{Source: arXiv:2605.05484v1 (CC BY 4.0)}
\begin{document}
\maketitle

\noindent\emph{Source and licence.} Multifractal analysis of power means for the Schneider map on $p\mathbb{Z}_{p}$. Version arXiv:2605.05484v1 (CC BY 4.0), distributed by arXiv under the Creative Commons Attribution 4.0 International licence, http://creativecommons.org/licenses/by/4.0/, as recorded in the arXiv licence metadata for this exact version and shown on the abstract page https://arxiv.org/abs/2605.05484v1. Full licence notice: \texttt{licenses/arxiv-2605.05484-CC-BY-4.0.txt}.\par\smallskip

\noindent\textit{Editorial scope.} Counted unit: the article's proposition PropCalculosHar (source0.tex lines 697--706). The article's complete original proof of it is delivered with it (source0.tex lines 709--729, one \texttt{proof} environment of 21 lines). No mathematics is written, repaired or re-derived: the statement and the proof are the article's own lines, in the article's own notation, and no retained line is substituted or re-flowed.  Retained uncounted as the source-local context the proof needs: lines 538--540; lines 542--544; lines 627--630.  Nothing was omitted from any retained span, so the package carries no omission record.  No retained line contains a citation, so the wrapper carries no bibliography and no bibliography entry.  No retained line contains a figure environment, an image-inclusion command or drawing code, so no image is delivered and the package declares no asset dependency.  Editorial formatting only, in addition: an AMS \texttt{amsart} preamble that restates the article's own declarations and macros used by the retained lines, namely the environments \texttt{proposition} on separate counters, together with the standard packages those lines need.  The retained environments print in this document as Proposition 1 in order of appearance, whereas the article prints the same items with the numbering of its own full document.  The complete native source of the article as distributed in the arXiv e-print for this exact version is bundled unchanged as \texttt{sources/199959/source0.tex}.
\begin{align}\label{QpowermeanP}
    M_{q}\left(\left\{a_1(x),...,a_n(x)\right\}\right)^{q}=\dfrac{1}{n(\log p)^{q}}\sum^{n-1}_{k=0}\left(\log \psi (T^{k}x)\right)^{q}.
\end{align}

\begin{align}\label{QpowermeanG}
    \log M_0(\{a_1(x),...,a_n(x)\})=-\log\log p+\dfrac{1}{n}\sum_{k=0}^{n-1} \log \log \psi (T^{k}x).
\end{align}

\begin{proposition}\label{derivative in zero}
The following identity holds:    
\[\left.\dfrac{d}{ds}\Li_s(1/p)\right|_{s=0}=-\sum_{k=1}^\infty \dfrac{\log k}{p^k}.\]
\end{proposition}

\begin{proposition}\label{PropCalculosHar}
For $\mu_p$-almost every $x\in p\Zp\setminus F$, the following equations hold:
    %\vspace{-1mm}
    \begin{equation}\label{Kqctp}
    M_{q}(x)=\left((p-1)\Li_{-q}(1/p)\right)^{\frac{1}{q}} \text{  for $q\neq 0$, and }
    \end{equation} %\vspace{-2mm}
    \begin{equation}\label{K0ctp}
    M_{0}(x)=\exp\left(-(p-1)\left.\dfrac{d}{ds}\Li_s(1/p)\right|_{s=0}\right)
    \end{equation}
\end{proposition}

\begin{proof}
    From equation \eqref{QpowermeanP} we have
    \[M_q(x)^q=\dfrac{1}{(\log p)^q}\lim_{n\to \infty}\dfrac{1}{n}\sum_{k=0}^{n-1}(\log \psi (T^ix))^q. \]
    By Birkhoff ergodic theorem, this last expression coincides for Haar-almost every $x$ with
    \begin{align}\label{IntKq}
        \dfrac{1}{(\log p)^q}\int_{p\Z_p} (\log \psi(x))^q d\mu_p.
    \end{align}
    On the other hand, this integral can be computed as 
    \begin{align*}
        \dfrac{1}{(\log p)^q}\int_{p\Z_p} (a_k\log p)^q d\mu_p&=\sum_{k\geq 1}k^q \mu_p \left(p^k\Z_p \setminus p^{k+1}\Z_p \right)\\
        &=(p-1)\sum_{k\geq 1}\dfrac{k^q}{p^k}\\
        &=(p-1)\Li_{-q}(1/p).
    \end{align*}
Equation \eqref{Kqctp} follows from equation \eqref{IntKq}. The proof of equation \eqref{K0ctp} follows the same lines as the previous one.

From equation \eqref{QpowermeanG}, and Birkhoff ergodic theorem we have for $\mu_p$-almost every $x\in p\Z_p\setminus F$
 \[\log M_0(x)=-\log \log p +\int_{p\Z_p} \log \log \psi(x)d\mu_p\]
 evaluating the integral, we obtain
 \[\log M_0(x)=(p-1)\sum_{k=1}^\infty \dfrac{\log k}{p^k}.\]
 We conclude the proof by proposition \ref{derivative in zero}.
\end{proof}

\end{document}
